\mysection{Introduction}

Serving LLMs is costly for providers such as Hugging Face and Alibaba Cloud, which must host thousands of base and fine-tuned models, many with low request volumes yet mandatory availability~\cite{aegaeon, slora, xiang2025}. As token prices fall and models vary widely in size and workload patterns, reducing inference cost while preserving performance has become a central objective. The core challenge, however, is not hardware expense but chronic under-utilization: to satisfy strict SLOs, industry practice often dedicates a \emph{model-parallel GPU group} (\ie, one or more GPUs jointly serving a single model instance), guaranteeing immediate responsiveness but wasting substantial capacity. This strategy is particularly inefficient for models with bursty or sparse traffic, and production analyses show that GPU duty cycles commonly fall below 30\%~\cite{aegaeon, xiang2025}.

To improve utilization, prior work has explored various GPU sharing mechanisms at different granularities. 
Because auto-regressive LLM inference is intrinsically memory-bound, tail latency is governed primarily by memory stalls (\eg, page faults, PCIe migrations). 
Moreover, compute sharing is already well studied\textemdash both via hardware mechanisms such as NVIDIA MIG~\cite{nvidiaMIG} and vGPU~\cite{nvidia-vgpu}, and via software solutions~\cite{nvidiaMPS,tally@asplos23,fu2024serverlessllm,lithos@sosp25,gpreempt@atc25,xsched@osdi25}. In contrast, effective memory sharing remains comparatively underexplored. We therefore center this work on advancing effective GPU memory sharing.

The work of memory sharing for LLM serving falls into two categories: space sharing and time sharing. 
Space sharing (\eg, MuxServe~\cite{muxserve}) colocates multiple models on the same GPU to utilize unused memory and compute cycles, which is effective for smaller, low-traffic models that are continuously accessed. Time sharing (\eg, QLM~\cite{qlm} and Aegaeon~\cite{aegaeon}) swaps models in and out of GPU memory, which is ideal for handling sporadic requests to models as the resources can be freed up when models are idle. While these strategies are effective for specific cases, our analysis of a broad range of production serving traces demonstrates that they fail to adapt to more complex and realistic usage scenarios. 

In particular, we have conducted a comprehensive analysis of production traces drawn from major inference providers (\ie, Hyperbolic and Novita AI) as well as independent evaluation platforms (\ie, Chatbot Arena). These traces encompass 11 days to 16 months of traffic across 58 heterogeneous models. Our results indicate that real-world workloads display diverse, rapidly changing patterns that elude any single, fixed sharing mechanism. 

\MyPara{Key observations.} \underline{At the model level}, deployed models show a \emph{bursty-group} behavior: when a burst occurs, requests concentrate on a particular subset of models, and that subset shifts over time, similar to an application’s dynamically changing working set. This pattern is driven by modern compound AI systems and agentic pipelines, which depend on a central reasoning LLM together with a number of fine-tuned or LoRA models for specialized agent tasks. \underline{At the request level}, detailed inspection of the traces confirms that requests are extremely \emph{dynamic and volatile} and fluctuate rapidly and unpredictably, limiting the value of existing GPU-sharing techniques. 
Space sharing, alone, fails to adapt to longer idle periods; for instance, in the Novita AI trace, where models are idle for over 70\% of the time, space sharing creates significant memory fragmentation, leaving 50\% of GPU memory occupied but unused. In contrast, time sharing, alone, struggles with rapid workload fluctuations. The significant latency overhead of swapping (\eg, several seconds) causes severe thrashing during traffic bursts, resulting in a spike of requests violating their latency SLOs due to queuing delays. 

Maximizing resource utilization while preserving SLO attainment requires integrating both strategies: space sharing for colocating low-traffic models and time sharing for burst handling and idle-resource reclamation. However, mingling these approaches is far from trivial due to the fundamentally different challenges they must address\textemdash the flexibility of memory balancing for space sharing vs. the efficiency of model swapping for time sharing. 
One could attempt to switch spatial and temporal sharing strategies dynamically;
yet, co-locating models with conflicting sharing needs creates resource contention. Specifically, a model requiring time-sharing typically relies on preempting resources to maximize burst performance, which directly contradicts the persistent resource availability required by a model colocated through space-sharing. Without a mechanism to fluidly mediate these conflicting operational modes, the system is forced into rigid configurations that cannot adapt to dynamic load changes.

\MyPara{Insight: a memory-centric view to unify space- and time-sharing.} 
Our central observation is that \emph{GPU memory is the unifying bottleneck, and each sharing strategy is concerned with a different portion of that memory}: time-sharing focuses on efficiently swapping in \emph{LLM weights}, whereas space-sharing determines how to scale \emph{KV-cache} capacity among concurrently served models. This perspective parallels \emph{memory ballooning} in virtualization, where a hypervisor dynamically reclaims memory from guest VMs and reallocates it where needed. In the context of LLM co-serving, an analogous “balloon driver” can reclaim memory from idle models to swap in the weights of new models (supporting time-sharing) or from models with low request rates by shrinking their KV-cache reservations (supporting space-sharing). Treating GPU memory as an elastic resource allows the system to fluidly shift between sharing modes and even enable both simultaneously (\ie, apply time-sharing to only a subset of colocated models), ensuring active models obtain the memory required to satisfy their SLOs while maximizing cluster-level efficiency.

\MyPara{\tool.} We developed \tool, a new memory-centric GPU sharing system for multi-LLM serving. At its core, \tool implements the concept of memory ballooning for LLMs to enable \emph{demand-aware, cross-model} memory sharing, a capability missing in existing systems but critical for unifying spatial and temporal strategies. \tool ensures that models can promptly expand their memory footprint to meet latency SLOs, while unused memory is efficiently harvested from others to maximize overall cluster utilization. The design of \tool addresses two key challenges.

\MyPara{Challenge 1: How to enable flexible memory ballooning for LLMs?} Today's LLM serving engines (\eg, SGLang and vLLM) are designed for single-model serving and adopt \emph{per-model static} memory allocation. While techniques like \codeIn{PagedAttention} manage memory efficiently within a model, they operate at the application level and cannot harvest memory across model boundaries.

Flexible memory ballooning requires a mechanism to dynamically resize the physical memory footprint of a model at runtime without disrupting its execution. \tool tackles this through a model balloon driver for serving engines (\S\ref{sec:sys}). 
The model balloon driver decouples virtual and physical GPU memory: serving engines reserve a large virtual address space during initialization, but physical memory is allocated and mapped only on demand. This design enables transparent, fine-grained memory redistribution at a 2\,MB granularity with millisecond-level overhead,
supporting diverse model architectures with zero code change to attention kernels.

\MyPara{Challenge 2: How to place models and schedule requests to maximize SLO attainment?} While the balloon driver provides the mechanism for elasticity, maximizing SLO attainment requires intelligent policies to decide \emph{when} and \emph{where} memory should be allocated. Without proper coordination, contention for shared memory can lead to thrashing and performance degradation.
\tool solves this through decoupled model placement and request scheduling. First, the fast-changing bursty groups (\S\ref{sec:motivation:model}) require a placement algorithm that can quickly decide where each model should run. \tool tackles this with a global demand-complementary placement algorithm that minimizes the KV pressure ratio (KVPR), balancing model memory demand with GPU capacity (\S\ref{sec:sch-global}).
Second, volatile request arrivals and diverse SLO targets require estimating each request’s time slack so the memory manager can steer memory toward latency-sensitive tasks. \tool achieves this via a local slack-aware arbitration algorithm, which leverages a GPU-level request queue to prioritize requests based on their slack (\S\ref{sec:sch-gpu}).

\MyPara{Results.} 
We implemented \tool on top of SGLang~\cite{sglang@arxiv23}, a widely used LLM serving engine. We evaluated \tool comprehensively with two production traces, and with \textbf{58 representative LLMs of varying architectures and sizes, on GPU clusters with up to 32 H100 GPUs. }
Our results show that \tool achieves up to 3.3$\times$ higher time-to-first-
token (TTFT) SLO attainment and $2\times$ higher time-per-output-token (TPOT) SLO attainment given the same number of GPUs.
When targeting the same level of SLO attainment, \tool achieves up to over 2$\times$ cost reduction or 3.5$\times$ more requests compared to the state-of-the-art. 

\tool's balloon driver, referred to as \codeIn{kvcached}, has been open-sourced and actively maintained at \url{https://github.com/ovg-project/kvcached}. \codeIn{kvcached} has been adopted by many industry users to support production workloads in clusters of more than 10K GPUs.
We report usage analyses from two companies at the end of \S\ref{sec:eval}. 
